\documentclass[11pt]{article}
\usepackage{fontspec}
\setmainfont{Times New Roman}
\setsansfont{Arial}
\setmonofont{Consolas}
\usepackage{amsmath,amssymb}
\usepackage{geometry}
\usepackage{microtype}
\geometry{a4paper,margin=3cm}
\setlength{\parindent}{1.25em}
\setlength{\parskip}{0pt}
\emergencystretch=30em
\allowdisplaybreaks
\raggedbottom
\providecommand{\mM}{\mathfrak M}
\providecommand{\mG}{\mathfrak G}
\providecommand{\mS}{\mathfrak S}
\providecommand{\mP}{\mathfrak P}
\providecommand{\mQ}{\mathfrak Q}
\providecommand{\mU}{\mathfrak U}
\providecommand{\ma}{\mathfrak a}
\providecommand{\mx}{\mathfrak x}

\begin{document}

% Unit: Paper 17 Section 06. Direct Interslavic Latin translation from the German final-audited source slice.

\setcounter{footnote}{14}
\subsection*{\S~6. Teoremy konečnosti.}

V tom paragrafu trěba pokazati, že vsaku grupu ostatkov možno predstaviti kako sumu konečno mnogih ireducibilnyh sostavnyh česti. Za to najprvo prenosimo Hilbertov teorem o modulnom bazisu.

Od nyně koristimo specijalnejše predpoloženja formulovane v \S~2: množenje \(\xi_1,\ldots,\xi_n\) medžu soboju jest komutativno, a zakon množenja polinomov ima formu
\begin{equation}
 \xi_i a=a_i\xi_i+b_i,
 \qquad (a_i\ne0)\quad(i=1,\ldots,n).
\tag{6}
\end{equation}
Dosta jest provesti dokaz samo za moduly. Bo ako \(\mS\) jest kaka-koli sistema polinomov, tvorimo iz nje proizhodny modul \(\mM\), dodavajuči k \(\mS\) vse polinomy, ktore možno sostaviti iz konečno mnogih polinomov \(G_1,\ldots,G_k\) iz \(\mS\) v formě \(A_1G_1+\cdots+A_kG_k\). Ako nyně \(F_1,\ldots,F_l\) jest modulny bazis od \(\mM\), i ako
\[
 F_i=\sum_j A_{ij}G_j\qquad(i=1,\ldots,l),
\]
togda konečno mnoge polinomy \(G_j\), ktore nastupajut v tih predstavjenjah, tvore bazis od \(\mS\).

Dokaz za modulny bazis provodimo po Gordaně (\emph{Göttinger Nachrichten} 1899). On razpada se na dvě česti. Prva odnosi se k sistemě \(\mP\) iz beskonečno mnogih produktov stepenev i ostavaje v sučnosti zachovana zaradi komutativnosti \(\xi_1,\ldots,\xi_n\). Produkty stepenev
\[
 P=\xi_1^{a_1}\cdots\xi_n^{a_n},
\]
iz ktoryh vsak zapisuje se v sistemu samo jedin raz, budut uredžene po rastučem stepenu i vnutri istogo stepena někojako opreděljeno. Razmatrjamo nyně podsistemu \(\mQ\) od \(\mP\), ktora sadrži vse take produkty stepenev \(P=Q\), ktore ne sut dělime nijednym predhodnym \(P\). Togda nijeden \(Q\) ne jest dělimy ni pozdnějšim \(Q\), bo ty sut ili vyšego stepena ili pri istom stepenu različne od ranjejših. Naprotiv vsaky \(P\) iz \(\mP\) jest dělimy najmene jednym \(Q\). Bo neha protivno \(P_m\) bude prvym \(P\), ktory ne jest dělimy nijednym \(Q\). Togda \(P_m\) ne može sam byti \(Q\), znači jest dělimy najmene jednym predhodnym \(P\); a poneže toj po predpoloženju jest dělimy někojim \(Q\), sam \(P_m\) jest takože dělimy \(Q\), što daje protivrěčje.

Tvrdimo nyně, že čislo \(Q\) jest konečno. V samoj stvari, neha
\[
 Q_0=\xi_1^{h_1}\cdots\xi_n^{h_n}
\]
bude prvy \(Q\). Poneže libovoljny \(Q_\lambda=\xi_1^{h_{1\lambda}}\cdots\xi_n^{h_{n\lambda}}\) ne jest dělimy \(Q_0\), za vsaky \(Q\) mora byti najmene jedno \(h_{\rho\lambda}<h_\rho\); neha \(\lambda\) vsaky raz označa najmanjši indeks, za ktory to važi. Sbiramo vse \(Q\), za ktore \(h_{\rho\lambda}\) ima tu samu fiksovanu vrědnost \(c_\lambda\), v jeden klas \(K(c_\lambda)\). Čislo tih klasov jest konečno, bo \(c_\lambda<h_\lambda\). Ale i v vsakom klasu ima samo konečno mnogo elementov. Bo ako \(Q\) iz klasa \(K(c_\lambda)\) položimo ravnym \(\xi_\lambda^{c_\lambda}Q'\), togda \(Q'\) takože tvore sistemu produktov stepenev, iz ktoryh nijeden ne jest dělimy drugim, i sadrže samo \(n-1\) proměnnyh. Pri \(n=2\) sadrže \(Q'\) tedy samo jednu proměnnu; jih čislo jest togda \(=1\), znači konečno, tako že čislo može byti vzęto kako konečno i za \(n-1\) proměnnyh. Tym tvrdženje jest dokazano.

Vtoro, neha \(\mM\) bude modul iz polinomov, ktoryh členy sut leksikografično uredžene. \(\mP\) bude sistema produktov stepenev, ktore nastupajut v najvyššem členu, vsak zapisany samo jedin raz. Togda sistemě \(\mQ=(Q_1,\ldots,Q_\rho)\) odpovědaje sistema konečno mnogih polinomov \(F_1,\ldots,F_\rho\), tako že vsakomu \(Q_i\) prirědžujemo někoj polinom iz \(\mM\), ktorego najvyšši člen jest ravny \(b_iQ_i\), takže možno položiti \(F_i=b_iQ_i+\cdots\) s nižšimi členami.

Neha nyně \(F=bP+\cdots\) bude libovoljny polinom iz \(\mM\), kde \(P=\xi_1^{\beta_1}\cdots\xi_n^{\beta_n}Q_i\). Tvorimo produkt
\[
 \xi_1^{\beta_1}\cdots\xi_n^{\beta_n}F_i
 =\xi_1^{\beta_1}\cdots\xi_n^{\beta_n}(b_iQ_i+\cdots)
 =c\,\xi_1^{\beta_1}\cdots\xi_n^{\beta_n}Q_i+\cdots
 =cP+\cdots.
\]
Tu \(c\ne0\) po (6); i vse slědujuče členy sut nižše členy leksikografičnogo reda, bo množenje, što se tiče eksponentov, jest komutativno do nižših členov.\footnote{Iz togo vidi se, že vměsto zakona (6) možno by bylo upotrěbiti i nekoliko menje ograničujuči zakon
\[
 \xi_i a=a_i\xi_i+a_{i,i+1}\xi_{i+1}+\cdots+a_{in}\xi_n+b_i\qquad(a_i\ne0)
\]
pri někojej fiksovanoj numeraciji \(\xi\).} Razlika
\[
 F-\frac bc\,\xi_1^{\beta_1}\cdots\xi_n^{\beta_n}F_i
\]
sadrži togda samo nižše členy i znovu naleži k \(\mM\). Postupok tedy prekida se po konečno mnogih krokih, tako že \(F_1,\ldots,F_\rho\) jest razpoznany kako modulny bazis.

\textbf{Teorem III.} \textit{Vsaka sistema \(\mS\) iz beskonečno mnogih polinomov ima konečny modulny bazis \(G_1,\ldots,G_k\), tako že vsaky polinom \(G\) iz \(\mS\) dopušča predstavjenje \(G=A_1G_1+\cdots+A_kG_k\).}

Iz togo neposrědnje slěduje

\textbf{Dodatok.} \textit{Vsaka sistema \(\mS\) iz beskonečno mnogih klasov ostatkov po modulu \(\mM\) ima bazis \(\mx_1,\ldots,\mx_k\), tako že vsaky \(\mx\) iz \(\mS\) jest predstavimy v formě \(\mx=A_1\mx_1+\cdots+A_k\mx_k\).}

Imenno, vsakomu klasu ostatkov prirědžujemo někoj reprezentant i razmatrjamo sistemu tih reprezentantov zajedno s vsimi polinomami iz \(\mM\).

Nyně prihodimo k dokazu slědujučego teorema.

\textbf{Teorem IV.} \textit{Vsaka grupa ostatkov može byti predstavljena kako suma konečno mnogih ireducibilnyh sostavnyh česti; tedy po Teoremu II vsaky modul može byti predstavjen kako najmene obče mnogokratno konečno mnogih ireducibilnyh modulov, ktore tvore totalno vzaimno prostu sistemu.}

Bo neha protivno \(\mG\) bude suma beskonečno mnogih sostavnyh česti \(\mU_1,\mU_2,\ldots\) (sravni \S~5), i neha \(\ma_1,\ma_2,\ldots\) budut naležeče jedinice, ktore po opreděljenju vse sut različne od \(0\). Togda tvorimo sistemu \(\mS\) vseh tih jedinic; po Dodatku k Teoremu III ona ima konečny bazis \(\ma_{i_1},\ldots,\ma_{i_v}\) s \(i_1<\cdots<i_v\). Neha nyně \(\mu>i_v\). Togda važi
\[
 \ma_\mu=\mx_1\ma_{i_1}+\cdots+\mx_v\ma_{i_v}.
\]
To jest linearna relacija medžu jedinicami, ktora po opreděljenju jest izključena. Pri vsakom additivnom razkladu grupy ostatkov zato po konečno mnogih krokih nastupa ireducibilna sostavna čest, ktoru zaradi komutativnosti razklada možemo postaviti na početok. Možemo tedy bez ograničenja obče važnosti same \(\mU_i\) smatrjati ireducibilnymi. Ale po gornjem zaključku može nastupiti samo konečno mnogo sostavnyh česti \(\mU_i\), i tym tvrdženje jest dokazano.

\end{document}
